
\[
0\longrightarrow
\operatorname{Gr}^{k-d}K^{\bullet}(Y)_{\mathbb Q}
\xrightarrow{u^{\gr}}
\operatorname{Gr}^{k-1}K^{\bullet}(Y')_{\mathbb Q}
\times
\operatorname{Gr}^{k}K^{\bullet}(X)_{\mathbb Q}
\xrightarrow{v^{\gr}}
\operatorname{Gr}^{k}K^{\bullet}(X')_{\mathbb Q}
\]
\uline{où on pose}, \uline{pour} \(y\in \operatorname{Gr}^{k-d}K^{\bullet}(Y)_{\mathbb Q}\) :
\[
  u^{\gr}(y)=
  \bigl(g^{\gr}(y).(-1)^{d-1}c^{d-1}(F),-i_{\gr}(y)\bigr),
\]
\uline{et pour} \(y'\in \operatorname{Gr}^{k-1}K^{\bullet}(Y')_{\mathbb Q}\) \uline{et} \(x\in \operatorname{Gr}^{k}K^{\bullet}(X)_{\mathbb Q}\) :
\[
  v^{\gr}(y',x)=j_{\gr}(y')+f^{\gr}(x),
\]
\uline{les notations étant celles de} 3.1 \uline{et} 3.3, \uline{et} \(c^{d-1}(F)\) \uline{désignant la} \((d-1)\)-\uline{ième classe de Chern de} \(F\) (V 6.7). \uline{En outre}, \(u^{\gr}\) \uline{admet un inverse à gauche} \(u'^{\gr}\), \uline{défini par} \(u'^{\gr}(y',x)=g_{\gr}(y')\).

\uline{Si de plus} \(X\) \uline{est noethérien et} \(Y\) \uline{régulier}, \uline{on a la suite exacte} :
\[
0\longrightarrow
\operatorname{Gr}^{k-d}K^{\bullet}(Y)_{\mathbb Q}
\xrightarrow{u^{\gr}}
\operatorname{Gr}^{k-1}K^{\bullet}(Y')_{\mathbb Q}
\times
\operatorname{Gr}^{k}K^{\bullet}(X)_{\mathbb Q}
\xrightarrow{v^{\gr}}
\operatorname{Gr}^{k}K^{\bullet}(X')_{\mathbb Q}
\longrightarrow0
\]

\emph{Démonstration.} Notons qu'on pourrait donner du cas \(Y\) régulier une démonstration purement algébrique, en termes de théorie des \(\lambda\)-anneaux, utilisant la description 4.7 de la situation. On va au lieu de cela calquer la démonstration de 3.7.

On sait que les homomorphismes « image inverse » \(f^*,g^*\), étant des \(\lambda\)-homomorphismes, sont compatibles aux filtrations de \(\lambda\)-anneaux, et par suite ils définissent des homomorphismes sur les gradués associés, d'où les homomorphismes :
\[
  f^{\gr}:\operatorname{Gr}^kK^{\bullet}(X)\to\operatorname{Gr}^kK^{\bullet}(X');
  \qquad
  g^{\gr}:\operatorname{Gr}^{k-d}K^{\bullet}(Y)\to\operatorname{Gr}^{k-d}K^{\bullet}(Y'),
\]
et les homomorphismes analogues obtenus après tensorisation par \(\mathbb Q\). Par ailleurs, 4.6 montre qu'après tensorisation par \(\mathbb Q\), les homomorphismes \(i_*,j_*\) passent aux gradués associés, et définissent
donc :
\[
  i_{\gr}:\operatorname{Gr}^{k-d}K^{\bullet}(Y)_{\mathbb Q}
  \to \operatorname{Gr}^{k}K^{\bullet}(X)_{\mathbb Q};
  \qquad
  j_{\gr}:\operatorname{Gr}^{k-1}K^{\bullet}(Y')_{\mathbb Q}
  \to \operatorname{Gr}^{k}K^{\bullet}(X')_{\mathbb Q}.
\]
